\chapter{Guidance Design}

\section{Mask}

The mask is a set of weigths (sum to 1) over a tensor having the same dimension as the weather states, non-zero at the variable and region of interest.

We experimented with a gaussian mask, but realized in our experiments that the best thing to do is equal weighting inside the region.
The correct way to define the region of interest is to choose a greater mask size instead of using some weighting scheme fading out of an epicenter. The mask size should be chosen to cover the entire area of interest, ensuring that all relevant data is included in the analysis. A larger mask size can help to reduce noise and improve the accuracy of the results, while a smaller mask size may exclude important information. It is important to carefully consider the appropriate mask size for each specific application to achieve optimal results.
Either way we (most probably) loose the global coherence for the guided weather states, so we might as well define directly the region of interest with a mask. 
Additionally, this enables us to actually make statements about "the region of interest",
having equal weights everywhere.
Introducing different weigths might bias the generation in a way that is not desired.



\section{Target trajectory}

In order to define plausible average trajectories of change, we define following procedure:
1. Select start state
2. Define mask
3. Retrieve historical states that are consistent with the start state and the mask (+- some tolerance)
4. Compute the distribution of average trajectories of change for the historical states
5. Filter a quantile of the distribution to define the target trajectory (e.g. 0.95 quantile for extreme events)
6. Define a gaussian distribution at each time step of the target trajectory, with mean equal to the target trajectory and variance equal to the variance of the historical states at that time step.
7. Sample enough trajectories from the gaussian distribution to define a set of target trajectories that will be used to guide the generative model.
8. Select the target trajectories of interest in terms of extremeness to build scenarios.
9. Save the trajectories and corresponding states used to define the sampling distribution. These will be used to evaluate the plausibility of the generated trajectories as means of comparison.
10. Run benchmark tests and evaluate.

This heuristic enables us to anchor the guidance target and define a benchmark at the same time.
More concretly, we avoid the issue of having to tune the strength of the guidance to achieve a certain plausible result.
We reverse the problem by defining the plausible result first and then designing an algorithm to achieve that result and checking that both the algorithm and its result are well behaved.

\section{Guidance loss}

The loss we use to guide the generative model is defined as follows:
...

One of the main contributions of this thesis is the realization that in the weather domain, we can define a guidance target, which can be achieved precisely.
By precisely I mean that the loss can be driven to zero.
While in the image domain, there is always a ground truth during training, at inference we play with different guidance strengths to identify a range where the generated images are good looking.
In the weather domain however, we can define a target that can be closed by tuning the guidance strength.

This is great because instead of tuning the guidance strength to find a range of plausible results, we can define what plausible results look like a priori (guidance target sampling trajectory defined through historical data and plausibility tests), and optimize the guidance strength to achieve this target. 



We essentially have a very similar setting to the guided diffusion for precipitation paper, but we make it more general in these ways:
\begin{itemize}
    \item we define a guidance target explictly, in contrast to observe results across different lambdas
    \item we experiment with different $a_t$ trajectories, in addition to guiding computing the gradient at the first step
    \item we test level of intensity by shifting the intenisty parameter from the flow (lambda) to the trajectory definition ($phi_n$)
    \item Our loss is driven to zero in terms of the difference from the target, instead of being minimized iteratively (a more explicit way to encode the target scenario)
\end{itemize}

\section{Guidance method}

Assumption about using the gradient.
We explicitly assume that in order to produce realistic result we have to tie the perturbations to the models' gradient.
We use the none-gradient tied perturbations as baseline for all our experiments.

The hope is also that the gradient is informative in the sense that:
if a method pushes a not secondary target variable to a certain distribution / average value, we would also expect that if we push this variable to the achieved average value, that the previously targeted variable would achieve the same result.
This we explicitly test in test5.

The base guidance algorithm we use works as follows:
...



\section{Guidance algorithms}

\subsection{CLOSE-GAP}

Greedy

Step is normalized to close the gap as prescribed by the guidance schedule.

\subsection{OPTIMIZE-GAIN}

Secant method

The step is allowed to take the gradient norm into consideration.
Since the loss is almost linear with respect to the guidance strength, we can use a secant method to find the guidance strength that minimizes the loss.

In this case we define the secant method as follows:


\subsection{OPTIMIZE-SCHEDULE}

The step is allowed to take the gradient norm into consideration.
The schedule is optimized to minimize the loss while also complying to a regularization term.
Possible regularization terms include the total guidance.
We can dream up arbitrary penalty terms such as angular defelction, ... .

The problem of this optimization is that the gradient applyed for the guidance at each step changes with respect to the schedule.
The optimization is therefore highly non-convex and very hard to solve.
In fact, the optimizer will converge very fast to a local minimum, highly influenced by the initialization of the schedule.

LBFGS